\begin{abstract}
A bank support desk must decide which waiting ticket an agent opens next, as high-severity issues such as fraud require immediate intervention under strict response windows. \af{Too many things cluttered in the opening sentence. Rephrase to clarify the problem}Published work on support ticket prioritization addresses \af{this as a classification problem which is trained with an offline dataset....and add how the rest is handled interms of}\sout{half of this operational decision by stopping at static classification and is mostly outside the banking domain}. This paper presents an operational triage and dynamic queue ordering framework that uses fine-tuned multilingual encoders to classify incoming tickets and compute dynamic urgency scores for helpdesk dispatch. Fine-tuned LaBSE encoders predict ticket intent, sentiment, and priority as continuous probability distributions, outperforming multilingual LLMs and classical baselines across English, Sinhala, Tamil, and their Romanized colloquial variants. \af{next sentence not clear...rephrase. Sentences do not start with because}Because intent accounts for 76.9\% of the uncertainty in priority, treating classification heads as independent evidence counts shared signal twice; an equal-weight log pool merges direct priority predictions with an intent-conditioned chain, reducing log loss by 27.4\% and resolving collinearity. We then specify a dynamic urgency scoring function that couples this combined posterior with elapsed waiting time, so that queue position reflects both predicted severity and accrued delay and no ticket can be starved indefinitely by construction; \af{dont add future work in abstract}its evaluation against dispatch policies under simulated load is left to future work. To train and evaluate the system, we expanded BANKING77 into a 5-track corpus of 65,385 tickets, pairing LLM-generated priority and sentiment labels, with sentiment-label quality assessed against a manually annotated gold subset.
\end{abstract}

